
\subsection{Méthode complémentaire: Facteurs de vue d'ordre 2}
\red	{ est-ce utile ? on verra plus tard}
La méthode des \textbf{facteurs de vue d'ordre 2} permet d'obtenir une approximation plus précise des facteurs de vue en prenant en compte la première interraction indirecte entre les surfaces. Cette amélioration est très utile lorsque les surfaces sont très grossièrement discrétisées. Elle repose sur l'introduction de facteurs de vue d'ordre 2, $\overline{J^2}$, qui prennent en compte la première interaction indirecte entre les surfaces et qui reposent sur les facteurs de vue d'ordre 1 :  $ \forall i, \quad J_i(\vec{r_i}) = J^1_i(\vec{r_i}) + \overline{J^2}_i - \overline{J^1_i} $. L'article \cite{Romain25} montre de la même manière que pour l'ordre 1 comment obtenir la formulation suivante :
\begin{equation}
    \overline{J^2_i} - \overline{J^1_i} - \rho_i \sum_{j=1}^{n} F_{ij}  \bigg(\overline{J^2_j} - \overline{J^1_j}\bigg) = \rho_i \sum_{j=1}^{n}  F_{ij} \epsilon_j \sigma \overline{T^4_j} + \rho_i \sum_{j = 1}^{n} \rho_j \sum_{k=1}^{n} F_{ijk} \overline{J^1_k} - \rho_i \sum_{j=1}^{n} F_{ij} \overline{J^1_j}
\end{equation}
avec $F_{ijk}$ le facteur de vue d'ordre 2, qui correspond à la fraction de radiation quittant la surface $i$ qui est intercepté par la surface $j$ après une première interraction avec la surface $k$. $F_{ijk} = \sum_{l=1}^{N} \frac{A_{i,l}}{A_i} \int_{A_{i,l}} \frac{1}{A_{i,h}} d^2r_i f_{ijk}(\vec{r_i})$. On construit enfin l'estimateur de Monte-Carlo de l'espérance pour les facteurs de vue d'ordre 2 $\forall \alpha = (h,l,m,l',m')$ :
%\begin{equation}
%    \hat{F}_{ijk}^{N_i, N_\theta, N_\varphi, N} = \sum_{h=1}^{N_i} \frac{A_{i,h}}{A_i} \sum_{l=1}^{N_\theta} \sum_{m=1}%^{N_\varphi} \alpha_{l,m} \sum_{l'=1}^{N_\theta} \sum_{m'=1}^{N_\varphi} \alpha_{l',m'} \hat{F}^{\alpha,N}_{ijk}
%\end{equation}

\begin{equation}
    \hat{F}_{ijk}^N = \frac{1}{N} \sum_{n=1}^N \hat{F}^n_{ijk}
\end{equation}
avec, toujours pour des poids unitaires sur les angles :
%\begin{equation}
%    \hat{F}_{ijk}^{\alpha,N} = \frac{1}{N} \sum_{n=1}^{N} 1_{A_j}\big(\vec{r}_i^{\alpha,n}, \vec{\Omega}_{ij}^{\alpha,n}\big) 1_{A_k} \big(\vec{r}_j^{\alpha,n}(\vec{r_i}, \vec{\Omega}_i), \vec{\Omega}_{jk}^{\alpha,n} \big)
%\end{equation} 
\begin{equation}
    \hat{F}_{ijk}^n =  1_{A_j}\big(\vec{r}_i^{\alpha,n}, \vec{\Omega}_{ij}^{\alpha,n}\big) 1_{A_k} \big(\vec{r}_j^{\alpha,n}(\vec{r_i}, \vec{\Omega}_i), \vec{\Omega}_{jk}^{\alpha,n} \big)
\end{equation} 

D'autres pistes d'amélioration peuvent être empruntées. L'objectif ici est de se concentrer sur les facteurs de vue d'ordre 1. Mais une liste non exhaustive de méthodes complémentaires est faite dans \cite{Romain25}. Elles ont pour but de compenser certaines approximations que peuvent engendrer les facteurs de vue. La méthode de Gebhart \cite{gebhart} permet notamment de corriger le biais introduit par l'hypothèse de radiosité uniforme. 



\subsection{Stratégie d'échantillonnage}
\red	{stade de brouillon, malheureusement non nécessaire pou l'instant}
\label{part:echanti}
\begin{equation}
    \hat{F_{ij}}^{N_i, N_\theta, N_\varphi, N} =  \sum_{h = 1}^{N_i} \frac{A_{i,h}}{A_i} \sum_{l = 1}^{N_{\theta}} \sum_{m = 1}^{N_{\varphi}} \alpha_{l,m} \hat{F}_{ij}^{h,l,m, N} \quad \text{avec } \hat{F}_{ij}^{h,l,m, N} = \frac{1}{N} \sum_{k = 1}^{N} 1_{A_j}(\vec{r_i}^{(h,l,m,k)}, \vec{\Omega_i}^{(h,l,m,k)}) 
    \label{eq:fv1}
\end{equation}

Un estimateur Monte-Carlo de la variance se construit de la même manière.


Le choix de partionner l'hémisphère en $N_\theta \times N_\varphi$ directions de lancement, motivé par \cite{partition} et mise en place dans \cite{Romain25} permet d'aboutir à une décomposition comme la suivante :


La tâche ici n'est pas d'étudier différentes constructions d'estimateurs. La fonction densité de probalilité choisie apparaît natuellement dans la définition de \ref{eq:MCfij}. On introduit donc $p_X = \frac{cos(X)sin(X)}{\pi}$, associée à la constante de normalisation $ \alpha_{l,m} $ conformément à \cite{Romain25}. Cet échantillonage uniforme sur la surface et sur les angles permet d'avoir une fonction poids associée unitaire. On fait donc $N$ tirages pour $(\vec{r_i}^{(l,m,k)}, \vec{\varphi_i}^{(l,m,k)}, \vec{\theta_i}^{(l,m,k)})$. L'estimateur Monte-Carlo de l'espérance :


des stratégies de stratifications seront discutées ultérièrement et donnent de meilleurs résultats de convergence qu'un échantillonage uniforme sur les deux angles \cite{Romain25}.
avec    $\alpha_{l,m}= \frac{\big(\varphi_{m+1}- \varphi_{m} \big)\big(sin^2(\theta_{l+1})-sin^2(\theta_{l})\big)}{2\pi}$.

\begin{equation}
    F_{ij} = \frac{1}{A_i} \int_{A_i} d^2\vec{r}_i \sum_{l = 1}^{N_{\theta}} \sum_{m = 1}^{N_{\varphi}} \int_{\theta_{l}}^{\theta_{l+1}} \int_{\varphi_{m}}^{\varphi_{m+1}} \bigg[\frac{cos(\theta_i) sin(\theta_i)}{\pi} 1_{A_j}(\vec{r_i}, \vec{\Omega_i})\bigg] d\varphi_i d\theta_i
\label{eq:MCfij}
\end{equation}




\subsection{Gabhart}
\red{ ça serait bien que j'ai le temps de l'implementer dans E2T, pour l'instant ça reste comme ça...}
\begin{equation}
B_{ij} = F_{ij} \epsilon_j + \sum_{k=1}^n F_{ik} (1 - \epsilon_k) B_{kj}
\end{equation}


\begin{equation}
Q_i = q_i A_i = A_i \epsilon_i \sigma T_i^4 - \sum_{j=1}^nA_j\epsilon_j \sigma T_j^4 B_{ij}
\end{equation}


La méthode des moindres carrés reformule la correction des facteurs de vue comme un problème d'optimisation sous contraintes. On cherche la matrice corrigée $F$ qui minimise l'écart quadratique pondéré à la matrice brute $F^0$ issue du calcul Monte Carlo :

\begin{equation}
    y = \sum_{i=1}^{n} \sum_{j=1}^{n} w_{ij} \left(F_{ij} - F_{ij}^0\right)^2
    \label{eq:moindre_carre}
\end{equation}

où les poids $w_{ij}$ permettent de pondérer chaque terme selon l'incertitude associée à son estimation Monte Carlo (les facteurs de vue estimés à partir d'un plus grand nombre de rayons pouvant par exemple se voir attribuer un poids plus élevé). Cette minimisation est soumise aux contraintes de fermeture et de réciprocité, qui s'écrivent sous forme d'égalités linéaires :

\begin{equation}
    \forall i, \quad \sum_{j=1}^{n} F_{ij} = 1 
    \label{eq:contrainte_fermeture}
\end{equation}

\begin{equation}
    \forall i \neq j, \quad A_j F_{ji} - A_i F_{ij} = 0
    \label{eq:contrainte_reciprocite}
\end{equation}

Une contrainte de positivité $F_{ij} \geq 0$ peut également être ajoutée si nécessaire, transformant le problème en un programme quadratique sous contrainte d'inégalité.

Plutôt que de résoudre ce problème par des méthodes générales de programmation non linéaire, \cite{renforcement95} propose une résolution reposant sur une interprétation géométrique du problème. Les facteurs de vue sont réorganisés sous la forme d'un vecteur colonne $f$ (obtenu en empilant les lignes de la matrice $F^0$), ce qui permet de réécrire les contraintes de fermeture et de réciprocité~\ref{eq:contraintes} sous la forme d'un système linéaire :

\begin{equation}
    R \cdot f = c
    \label{eq:systeme_contraintes}
\end{equation}

La solution de ce système peut être décomposée en deux composantes orthogonales : une solution particulière $f_{row}$, appartenant à l'espace ligne de $R$ et de norme minimale, donnée par

\begin{equation}
    f_{row} = R^T (R R^T)^{-1} c
    \label{eq:solution_particuliere}
\end{equation}

et une composante homogène appartenant au noyau de $R$, qui s'exprime comme une combinaison linéaire des vecteurs d'une base $N_b$ de ce noyau :

\begin{equation}
    f - f_{row} = N_b x
    \label{eq:solution_homogene}
\end{equation}

Les facteurs de vue bruts $f^0$ ne vérifiant pas exactement les contraintes, l'équation~\eqref{eq:solution_homogene} n'admet pas de solution exacte : le vecteur de poids $x$ est alors déterminé au sens des moindres carrés, en projetant l'écart $(f^0 - f_{row})$ sur le noyau de $R$ :

\begin{equation}
    x = (N_b^T N_b)^{-1} N_b^T (f^0 - f_{row})
    \label{eq:solution_x}
\end{equation}

Le vecteur des facteurs de vue corrigés $f$, reconstruit à partir des équations~\eqref{eq:solution_particuliere}, \eqref{eq:solution_homogene} et \eqref{eq:solution_x}, constitue alors la projection de $f^0$ sur l'espace des solutions vérifiant exactement les contraintes de fermeture et de réciprocité. Contrairement à l'algorithme naïf, cette approche traite l'ensemble des $n^2$ facteurs de vue de façon symétrique et cohérente, sans privilégier arbitrairement une partie de la matrice (comme le fait l'algorithme naïf en ne conservant que les termes au-dessus de la diagonale) ni recourir à un processus itératif de convergence incertaine (comme l'algorithme de Leersum). Cette rigueur a cependant un coût : le vecteur $f$ comportant $n^2$ composantes, la matrice $N_b^T N_b$ à inverser dans l'équation~\eqref{eq:solution_x} est de taille de l'ordre de $n^2 \times n^2$, ce qui conduit à une complexité de calcul évoluant en $\mathcal{O}(n^4)$. Cette méthode est ainsi vraisemblablement inadaptée à des cas 3D de grande taille, pour lesquels le nombre de surfaces $n$ peut être élevé.





\section{taleau comparaitf}
\begin{table}[ht]
\centering

\renewcommand{\arraystretch}{1.}
\setlength{\tabcolsep}{3pt}
\scriptsize
\begin{tabular}{p{2.2cm}cccccc}

\textbf{Criterion} &
\textbf{Integration} &
\textbf{Unit Sphere} &
\textbf{Cross-String} &
\textcolor{codegreen}{\textbf{Monte Carlo}} &
\textbf{Ray Casting} &
\textbf{Matrix} \\

\hline
\hline

Computational cost&
Low&
Medium&
Low&
High&
\textcolor{red}{Very high}&
Low\\
\hline

Ray-tracing required &
No&
Yes&
No&
Yes&
Yes&
No\\
\hline

Required expertise &
Integration&
Integration&
Geometry&
Probability&
Programming&
Linear algebra\\
\hline

Analytical formulation &
complex&
Simple&
Simple&
Very complex&
Complex& 
\textcolor{red}{Simple}\\
\hline

Applicable to multiple surfaces &
\textcolor{red}{No}&
\textcolor{red}{No}&
\textcolor{red}{No}&
Yes&
Yes&
Yes\\
\hline

Applicable to closed enclosures &
No&
No&
No&
Yes&
\textcolor{red}{No}&
Yes\\
\hline

Typical geometries &
Simple&
Simple&
2D / curved&
Ours&
General&
Closed enclosures\\
\hline

\end{tabular}\\
\caption{Comparison of common methods for view factor computation. Adapted from \cite{methods_fv}}
\label{tab:methods}
\end{table}